\begin{table}[h]
\centering
\caption{Marginal Effects of Residential Zoning and Majority-Black Status on Highway Placement - Outside Highway Corridor}
\label{tab:marginal_effects_indir}
\begin{threeparttable}
\begin{tabular*}{0.6\textwidth}{@{\extracolsep{\fill}}l*{1}{r}}
\toprule
 & -5.044192 \\
\midrule
White Non-Residential & \makecell[tr]{0.032 \\ (0.009)} \\
White Residential & \makecell[tr]{0.004 \\ (0.002)} \\
Black Non-Residential & \makecell[tr]{0.008 \\ (0.004)} \\
Black Residential & \makecell[tr]{0.002 \\ (0.003)} \\
White Protection effect & \makecell[tr]{0.028^{***} \\ (0.010)} \\
Black Protection effect & \makecell[tr]{0.007 \\ (0.004)} \\
Disparate Protection (Black Protection - White Protection) & \makecell[tr]{-0.021^{***} \\ (0.008)} \\
\bottomrule
\end{tabular*}
\begin{tablenotes}[flushleft]
\footnotesize
\item This table contains predicted outcomes for each square in the sample based on their residential and majority-Black status, holding all other variables at their mean. The marginal effect of residential zoning and majority-Black status is calculated as the difference in predicted outcomes between squares with and without these characteristics. These marginal effects are calculated based on the coefficients reported in Table \ref{tab:indirect_results}. * p<0.10, ** p<0.05, *** p<0.01
\end{tablenotes}
\end{threeparttable}
\end{table}
